\clearpage
\section{정다각형과 원분체}
\label{sec:construction-regular-polygons}

\begin{learninggoal}
정$n$각형의 작도가능성을 원시 $n$차 단위근과 원분체의 차수로 판정한다. 가우스--방첼 정리를 증명하고 페르마 소수와의 관계를 설명한다.
\end{learninggoal}

단위원에 내접한 정$n$각형의 꼭짓점은
\[
  1,\zeta_n,\zeta_n^2,\dots,\zeta_n^{n-1},
  \qquad
  \zeta_n=e^{2\pi i/n}
\]
이다. 따라서 정$n$각형의 작도문제는 원시 $n$차 단위근 $\zeta_n$의 작도가능성 문제이다.

\begin{proposition}
$n\ge3$에 대해 다음은 동치이다.
\begin{enumerate}[label=(\roman*)]
  \item 정$n$각형을 자와 컴퍼스로 작도할 수 있다.
  \item $\zeta_n$이 작도가능하다.
  \item $\cos(2\pi/n)$과 $\sin(2\pi/n)$이 모두 작도가능하다.
\end{enumerate}
\end{proposition}

\begin{proof}
정$n$각형을 작도하면 인접한 두 꼭짓점 중 하나를 $1$로 놓았을 때 다음 꼭짓점이 $\zeta_n$이다. 역으로 $\zeta_n$을 작도하면 그 거듭제곱을 복소수 곱셈으로 차례로 작도하여 모든 꼭짓점을 얻는다. 마지막 두 조건은
\[
  \zeta_n=\cos(2\pi/n)+i\sin(2\pi/n)
\]
과 실수부·허수부의 작도가능성에서 동치이다.
\end{proof}

\subsection{오일러 함수 판정}

\begin{theorem}[정다각형 작도 판정]
\label{thm:regular-polygon-phi}
$n\ge3$에 대해 다음은 동치이다.
\begin{enumerate}[label=(\roman*)]
  \item 정$n$각형은 작도가능하다.
  \item $\varphi(n)$은 $2$의 거듭제곱이다.
  \item 원분체 $\Q(\zeta_n)/\Q$의 갈루아군은 $2$-군이다.
\end{enumerate}
\end{theorem}

\begin{proof}
원분체는 갈루아 확장이고
\[
  [\Q(\zeta_n):\Q]=\varphi(n),
  \qquad
  \Gal(\Q(\zeta_n)/\Q)\cong(\Z/n\Z)^\times.
\]
$\zeta_n$의 최소다항식은 원분다항식 $\Phi_n$이고 그 분해체가 바로 $\Q(\zeta_n)$이다. 따라서 최소다항식의 분해체 판정에 의해 $\zeta_n$이 작도가능할 필요충분조건은 이 분해체의 차수가 $2$의 거듭제곱인 것이다.
\end{proof}

\begin{example}
\[
\begin{array}{c|rrrrrrrrrrrrrrr}
 n&3&4&5&6&7&8&9&10&11&12&13&14&15&16&17\\
\hline
 \varphi(n)&2&2&4&2&6&4&6&4&10&4&12&6&8&8&16
\end{array}
\]
따라서 정$3,4,5,6,8,10,12,15,16,17$각형은 작도가능하지만, 정$7,9,11,13,14$각형은 작도가능하지 않다.
\end{example}

\subsection{가우스--방첼 정리}

\begin{definition}[페르마 수와 페르마 소수]
$r\ge0$에 대해
\[
  F_r=2^{2^r}+1
\]
을 $r$번째 페르마 수라 한다. $F_r$가 소수이면 이를 페르마 소수라 한다. 예를 들어
\[
  3,\ 5,\ 17,\ 257,\ 65537
\]
은 페르마 소수이다.
\end{definition}

\begin{lemma}
홀수 소수 $p$에 대해 다음은 동치이다.
\begin{enumerate}[label=(\roman*)]
  \item $p-1$은 $2$의 거듭제곱이다.
  \item $p$는 페르마 소수이다.
\end{enumerate}
\end{lemma}

\begin{proof}
(ii)$\Rightarrow$(i)는 정의에서 자명하다. 역으로 $p-1=2^m$이라 하자. $m=2^r u$로 쓰되 $u$는 홀수라 하자. 만약 $u>1$이면
\[
  p=(2^{2^r})^u+1
\]
이고 홀수 $u$에 대해 $a^u+1$은 $a+1$로 나누어지므로 $p$가 합성수가 된다. 따라서 $u=1$이고
\[
  p=2^{2^r}+1
\]
이다.
\end{proof}

\begin{theorem}[가우스--방첼 정리]
\label{thm:gauss-wantzel}
$n\ge3$에 대해 정$n$각형이 작도가능할 필요충분조건은
\[
  n=2^m p_1p_2\cdots p_r
\]
로 쓸 수 있는 것이다. 여기서 $m\ge0$이고 $p_1,\dots,p_r$은 서로 다른 페르마 소수이다.
\end{theorem}

\begin{proof}
소인수분해를
\[
  n=2^m\prod_{j=1}^r q_j^{e_j}
\]
로 쓰자. 오일러 함수의 곱셈성과 소수거듭제곱 공식에서
\[
  \varphi(n)
  =\varphi(2^m)\prod_{j=1}^r q_j^{e_j-1}(q_j-1)
\]
이다. 이 값이 $2$의 거듭제곱이면 홀수인 $q_j$가 인수로 나타날 수 없으므로 $e_j=1$이어야 한다. 또한 각 $q_j-1$이 $2$의 거듭제곱이어야 하므로 앞의 보조정리에 의해 $q_j$는 페르마 소수이다.

반대로 $n=2^m p_1\cdots p_r$이고 $p_j=2^{2^{a_j}}+1$인 서로 다른 페르마 소수라면
\[
  \varphi(n)
  =\varphi(2^m)\prod_{j=1}^r(p_j-1)
\]
은 $2$의 거듭제곱이다. 정리 \ref{thm:regular-polygon-phi}에 의해 정$n$각형은 작도가능하다.
\end{proof}

\begin{example}
\begin{itemize}
  \item $15=3\cdot5$이므로 정$15$각형은 작도가능하다.
  \item $17$은 페르마 소수이므로 정$17$각형은 작도가능하다.
  \item $51=3\cdot17$과 $85=5\cdot17$도 조건을 만족한다.
  \item $9=3^2$은 홀수 페르마 소수의 제곱을 포함하므로 정$9$각형은 작도가능하지 않다.
  \item $21=3\cdot7$은 $7$이 페르마 소수가 아니므로 정$21$각형은 작도가능하지 않다.
\end{itemize}
\end{example}

\subsection{갈루아 대응이 실제 작도를 보장하는 이유}

$\varphi(n)=2^s$이면
\[
  G=\Gal(\Q(\zeta_n)/\Q)
\]
는 위수 $2^s$인 아벨 $2$-군이다. 따라서
\[
  G=G_0\supset G_1\supset\cdots\supset G_s=1,
  \qquad [G_j:G_{j+1}]=2
\]
인 부분군열을 택할 수 있다. 갈루아 대응은 이를
\[
  \Q=K_0\subset K_1\subset\cdots\subset K_s=\Q(\zeta_n),
  \qquad [K_{j+1}:K_j]=2
\]
인 이차확장 탑으로 바꾼다. 각 단계가 하나의 제곱근 첨가로 표현되므로 이 정리는 단지 존재론적 판정에 그치지 않고, 원칙적으로 중첩 제곱근을 통한 작도 절차를 제공한다.

\begin{example}[정$17$각형]
\[
  \Gal(\Q(\zeta_{17})/\Q)
  \cong(\Z/17\Z)^\times\cong C_{16}.
\]
순환군 $C_{16}$에는
\[
  C_{16}\supset C_8\supset C_4\supset C_2\supset1
\]
이라는 지수 $2$의 부분군열이 있다. 이에 대응하는 네 번의 이차확장 탑이 $\zeta_{17}$을 포함하므로 정$17$각형은 작도가능하다. 실제 작도식이 길어지는 이유는 이 네 개의 제곱근 단계를 구체적인 좌표식으로 풀어 써야 하기 때문이다.
\end{example}

\begin{pitfall}
$\varphi(n)$이 단지 짝수라는 사실은 아무 의미가 없다. $n>2$이면 $\varphi(n)$은 언제나 짝수이다. 필요한 조건은 $\varphi(n)$의 소인수가 오직 $2$뿐이라는 것이다.
\end{pitfall}

\begin{checkpoint}
\begin{enumerate}[label=(\alph*)]
  \item 정$20$각형과 정$24$각형의 작도가능성을 각각 판정하라.
  \item 정$p$각형이 작도가능한 홀수 소수 $p$는 왜 페르마 소수여야 하는가?
  \item 정$17$각형의 작도에서 갈루아군의 부분군열은 어떤 역할을 하는가?
\end{enumerate}
\end{checkpoint}
